(1) 一维稳态傅里叶定律为 $j=-k(T)dT/dx$，没有内部热源则 $j$ 与 $x$ 无关。积分得 $jL=\int_{T_c}^{T_h}k(T)dT=k_0[(T_h-T_c)+\alpha(T_h^2-T_c^2)/2]$，故 $j=k_0[(T_h-T_c)+\alpha(T_h^2-T_c^2)/2]/L>0$。
(2) 从热面到位置 $x$，$\int_{T(x)}^{T_h} k(T)dT=jx$。记 $F(T)=T+\alpha T^2/2$，则 $F(T(x))=F(T_h)-(jx/k_0)$，并选使 $T_c\le T(x)\le T_h$ 的连续根；因 $F'(T)=1+\alpha T>0$，解唯一。
(3) 体熵产生密度是 $k(T)(dT/dx)^2/T^2$，因此总速率 $\dot S_{\rm gen}=A\int_0^L k(T)(T_x)^2/T^2\,dx=Aj\int_{T_c}^{T_h}dT/T^2=Aj(1/T_c-1/T_h)>0$。这也等于冷面熵流出 $Aj/T_c$ 减热面熵流入 $Aj/T_h$。
(4) 当 $\alpha\to0$，$j\to k_0(T_h-T_c)/L$、$T(x)\to T_h-(T_h-T_c)x/L$。保持两端温度不变而增大 $\alpha$ 时，若温度为正，整个积分 $\int k(T)dT$ 增大，故热流增加；不能把 $k(T)$ 随意用某一点温度替换而声称精确。
